% ==============================================================================
% PLANTILLA DE PRESENTACION - FILOSOFIA DEL DERECHO / UnADM
% ==============================================================================
% Uso rapido:
% 1) Duplica este archivo para una actividad concreta cuando necesites exponer.
% 2) Actualiza los datos editables: titulo, actividad, semana, docente y fecha.
% 3) Lee la planeacion correspondiente y traduce la consigna a una narrativa:
%    problema -> conceptos -> sistema juridico -> postura -> consecuencia.
% 4) Usa los modulos de esta plantilla segun el producto pedido:
%    mapa conceptual, cuadro comparativo, glosario, estudio de caso, reflexion,
%    interpretacion juridica, analisis constitucional o caso filosofico-juridico.
% 5) No pegues instrucciones de la planeacion. La diapositiva debe mostrar
%    criterio, no relleno.
% 6) Compilar desde la raiz del repositorio:
%    .\scripts\latexmk-build.ps1 .\UnADM\filosofia-del-derecho\presentacion-filosofia-del-derecho.tex
%
% Contrato de contenido heredado de los reportes de la materia:
% - Cada slide debe tener una pregunta o afirmacion principal.
% - Las fuentes sostienen una funcion analitica: definir, contrastar, aplicar,
%   validar, problematizar o justificar una postura.
% - Toda exposicion debe cerrar con postura personal argumentada.
% - Las actividades visuales se construyen como evidencia: mapa, cuadro, tabla,
%   matriz, linea de tiempo o diagrama, no como decoracion.
% ==============================================================================

\documentclass[
  spanish,
  aspectratio=169,
  xcolor={dvipsnames,table}
]{beamer}

% Lienzo ampliado 16:9: mantiene proporcion panoramica y da mas espacio util
% que el tamano por defecto de beamer.
\geometry{paperwidth=19.2cm,paperheight=10.8cm}

% Idioma, tipografia y recursos visuales
\usepackage[utf8]{inputenc}
\usepackage[T1]{fontenc}
\usepackage[spanish,es-tabla]{babel}
\usepackage[scaled=.96]{helvet}
\renewcommand{\familydefault}{\sfdefault}
\usepackage{array}
\usepackage{booktabs}
\usepackage{graphicx}
\usepackage{ragged2e}
\usepackage{tabularx}
\usepackage{tikz}
\usepackage{hyperref}
\usetikzlibrary{arrows.meta,calc,fit,matrix,positioning,shadows.blur,shapes.geometric}

\newcolumntype{Y}{>{\RaggedRight\arraybackslash}X}

% ------------------------------------------------------------------------------
% DATOS EDITABLES
% ------------------------------------------------------------------------------
\newcommand{\studentname}{Martin Jonathan de la Cruz}
\newcommand{\studentshort}{M. J. de la Cruz}
\newcommand{\studentid}{ES2611202040}
\newcommand{\studentemail}{martin.delacruz@unadmexico.mx}
\newcommand{\universityname}{Universidad Abierta y a Distancia de Mexico}
\newcommand{\facultyname}{Licenciatura en Derecho}
\newcommand{\coursename}{Filosofia del Derecho}
\newcommand{\coursecode}{030}
\newcommand{\activitytitle}{Plantilla de presentacion de Filosofia del Derecho}
\newcommand{\activitysubtitle}{Base para actividades existentes y nuevas}
\newcommand{\activityweek}{Semanas 2--8}
\newcommand{\teachingfigure}{Reyna Patricia Gonzalez Rodriguez}
\newcommand{\deliverydate}{\today}
\newcommand{\locationname}{Roma Norte, Ciudad de Mexico}
\newcommand{\generalbib}{bibliografia-unadm.bib}
\newcommand{\coursebib}{filosofia-del-derecho.bib}
\newcommand{\reporttemplate}{reporte-filosofia-del-derecho.tex}
\newcommand{\departmentlogo}{img/departamentos/UnADM.pdf}
\newcommand{\campusurl}{https://www.unadmexico.mx/}

% ------------------------------------------------------------------------------
% PALETA VISUAL
% Sobria e institucional, con verde UnADM y acento dorado.
% ------------------------------------------------------------------------------
\definecolor{unadmGreenDark}{HTML}{174A3A}
\definecolor{unadmGreen}{HTML}{5F8F3A}
\definecolor{unadmGold}{HTML}{B88A2A}
\definecolor{unadmPaper}{HTML}{F6F7F2}
\definecolor{unadmInk}{HTML}{1F2A24}
\definecolor{unadmMuted}{HTML}{5D6B66}
\definecolor{unadmSoft}{HTML}{E9EEE6}

\mode<presentation>{
  \usetheme{default}
  \usefonttheme{professionalfonts}
  \setbeamertemplate{navigation symbols}{}
  \setbeamertemplate{blocks}[rounded][shadow=false]
  \setbeamercovered{invisible}
}
\setbeamersize{text margin left=0.55cm,text margin right=0.55cm}

\hypersetup{
  colorlinks=true,
  urlcolor=unadmGreenDark,
  linkcolor=unadmGreenDark
}

\setbeamercolor{background canvas}{bg=white}
\setbeamercolor{normal text}{fg=unadmInk,bg=white}
\setbeamercolor{structure}{fg=unadmGreenDark}
\setbeamercolor{frametitle}{fg=white,bg=unadmGreenDark}
\setbeamercolor{title}{fg=white,bg=unadmGreenDark}
\setbeamercolor{subtitle}{fg=white}
\setbeamercolor{block title}{fg=white,bg=unadmGreen}
\setbeamercolor{block body}{fg=unadmInk,bg=unadmPaper}
\setbeamercolor{alerted text}{fg=unadmGold}
\setbeamerfont{title}{size=\LARGE,series=\bfseries}
\setbeamerfont{frametitle}{size=\large,series=\bfseries}
\setbeamerfont{block title}{series=\bfseries}
\setbeamertemplate{itemize item}{\textcolor{unadmGreen}{\large$\blacktriangleright$}}
\setbeamertemplate{itemize subitem}{\textcolor{unadmGold}{\scriptsize$\blacksquare$}}
\setbeamertemplate{enumerate item}{\textcolor{unadmGreenDark}{\insertenumlabel.}}

\setbeamertemplate{frametitle}{
  \nointerlineskip
  \begin{beamercolorbox}[wd=\textwidth,ht=1.02cm,dp=0.22cm,leftskip=0.35cm,rightskip=0.25cm]{frametitle}
    \begin{minipage}[c]{0.76\textwidth}
      \insertframetitle
    \end{minipage}
    \hfill
    \raisebox{-0.1cm}{\includegraphics[height=0.60cm]{\departmentlogo}}
  \end{beamercolorbox}
  \vspace{-0.04cm}
  {\color{unadmGold}\rule{\textwidth}{1.3pt}}
}

\setbeamertemplate{footline}{
  \leavevmode%
  \hbox{%
    \begin{beamercolorbox}[wd=.34\paperwidth,ht=2.7ex,dp=1ex,leftskip=1em]{author in head/foot}%
      \color{white}\studentshort
    \end{beamercolorbox}%
    \begin{beamercolorbox}[wd=.40\paperwidth,ht=2.7ex,dp=1ex,center]{title in head/foot}%
      \color{white}\coursename
    \end{beamercolorbox}%
    \begin{beamercolorbox}[wd=.26\paperwidth,ht=2.7ex,dp=1ex,rightskip=1em plus 1fil]{date in head/foot}%
      \color{white}\coursecode\hfill\insertframenumber/\inserttotalframenumber
    \end{beamercolorbox}%
  }%
  \vskip0pt%
}
\setbeamercolor{author in head/foot}{bg=unadmGreenDark}
\setbeamercolor{title in head/foot}{bg=unadmGreen}
\setbeamercolor{date in head/foot}{bg=unadmGreenDark}

\newcommand{\accentbar}{%
  \begin{tikzpicture}[remember picture,overlay]
    \path[use as bounding box] (0,0) rectangle (0,0);
    \fill[unadmGreen] (current page.south west) rectangle ([xshift=0.56\paperwidth,yshift=0.08cm]current page.south west);
    \fill[unadmGold] ([xshift=0.56\paperwidth]current page.south west) rectangle ([xshift=\paperwidth,yshift=0.08cm]current page.south west);
  \end{tikzpicture}
}

\newcommand{\sectionframe}[1]{%
  \begin{frame}[plain]
    \begin{tikzpicture}[remember picture,overlay]
      \path[use as bounding box] (0,0) rectangle (0,0);
      \fill[unadmGreenDark] (current page.south west) rectangle (current page.north east);
      \node[opacity=0.09,anchor=east] at ([xshift=-0.45cm]current page.east) {\includegraphics[width=0.43\paperwidth]{\departmentlogo}};
      \fill[unadmGreen] ([xshift=1.05cm,yshift=2.25cm]current page.south west) rectangle ([xshift=1.25cm,yshift=6.40cm]current page.south west);
      \fill[unadmGold] ([xshift=1.48cm,yshift=2.25cm]current page.south west) rectangle ([xshift=1.68cm,yshift=5.25cm]current page.south west);
      \node[
        anchor=west,
        align=left,
        text=white,
        text width=0.58\paperwidth,
        font=\fontsize{26}{31}\selectfont\bfseries
      ] at ([xshift=2.08cm,yshift=0.55cm]current page.west) {#1};
      \node[
        anchor=west,
        align=left,
        text=white!78,
        font=\large
      ] at ([xshift=2.08cm,yshift=-0.55cm]current page.west) {\coursename};
    \end{tikzpicture}
  \end{frame}
}

\newcommand{\unadmsection}[1]{%
  \section{#1}
  \sectionframe{#1}
}

\newcommand{\tagbox}[2]{%
  \begin{tikzpicture}[baseline=(n.base)]
    \node[
      rounded corners=2pt,
      fill=#1,
      text=white,
      inner xsep=5pt,
      inner ysep=2pt,
      font=\scriptsize\bfseries
    ] (n) {#2};
  \end{tikzpicture}%
}

\title[\coursecode]{\activitytitle}
\subtitle{\activitysubtitle}
\author[\studentshort]{\texorpdfstring{\studentname\\\footnotesize Matricula: \studentid\\\href{mailto:\studentemail}{\studentemail}}{\studentname, Matricula: \studentid}}
\institute[UnADM]{\universityname\\\facultyname\\\coursecode\ \textperiodcentered\ \coursename}
\date[\activityweek]{\locationname\\\activityweek\\\deliverydate}

\begin{document}

% Portada institucional
\begin{frame}[plain]
  \begin{tikzpicture}[remember picture,overlay]
    \path[use as bounding box] (0,0) rectangle (0,0);
    \fill[unadmGreenDark,opacity=0.96] (current page.south west) rectangle ([xshift=0.58\paperwidth]current page.north west);
    \fill[unadmGreen,opacity=0.92] ([xshift=0.58\paperwidth]current page.south west) rectangle (current page.north east);
    \node[anchor=center,opacity=0.10] at (current page.center) {\includegraphics[width=0.72\paperwidth]{\departmentlogo}};
    \node[anchor=north west] at ([xshift=0.58cm,yshift=-0.46cm]current page.north west) {\includegraphics[width=2.55cm]{\departmentlogo}};
    \fill[unadmGold] ([xshift=0.60\paperwidth,yshift=1.18cm]current page.south west) rectangle ([xshift=0.93\paperwidth,yshift=1.30cm]current page.south west);
    \fill[white,opacity=0.15] ([xshift=0.60\paperwidth,yshift=0.92cm]current page.south west) rectangle ([xshift=0.84\paperwidth,yshift=1.02cm]current page.south west);
  \end{tikzpicture}

  \vspace*{1.80cm}
  \begin{minipage}{0.55\paperwidth}
    {\color{white}\fontsize{24}{29}\selectfont\bfseries\raggedright \activitytitle\par}
    \vspace{0.22cm}
    {\color{white!86}\large \activitysubtitle}\\[0.32cm]
    {\color{unadmGold}\rule{0.82\linewidth}{1.3pt}}\\[0.42cm]
    {\color{white}\studentname}\\
    {\color{white!82}\footnotesize Matricula: \studentid}
  \end{minipage}
  \hfill
  \begin{minipage}{0.34\paperwidth}
    \raggedleft
    {\color{white}\Large\bfseries UnADM}\\[0.18cm]
    {\color{white!84}\footnotesize \universityname}\\[0.45cm]
    {\color{white}\small \facultyname}\\
    {\color{white!78}\small \coursename}\\[0.45cm]
    {\color{white!74}\footnotesize \coursecode\ \textperiodcentered\ \activityweek}\\
    {\color{white!74}\footnotesize \locationname}\\
    {\color{white!74}\footnotesize \href{\campusurl}{\campusurl}}
  \end{minipage}
\end{frame}

\begin{frame}{Uso de esta plantilla}
  \begin{columns}[T]
    \begin{column}{0.48\textwidth}
      \begin{block}{Antes de construir slides}
        \begin{itemize}
          \item Identifica el producto pedido por la planeacion.
          \item Extrae del reporte: objetivo, problema, conceptos, caso y postura.
          \item Selecciona solo los ejes que sostienen la exposicion.
          \item Sustituye los textos guia por contenido propio de la actividad.
        \end{itemize}
      \end{block}
    \end{column}
    \begin{column}{0.48\textwidth}
      \begin{block}{Regla de oro}
        La presentacion no es el reporte reducido: es la defensa visual de una
        postura. Cada slide debe mostrar una decision: que se afirma, con que
        concepto se sostiene y que consecuencia juridica produce.
      \end{block}
    \end{column}
  \end{columns}
  \accentbar
\end{frame}

\begin{frame}{Contenido}
  \tableofcontents
\end{frame}

\unadmsection{Contrato de la materia}

\begin{frame}{Ficha academica reutilizable}
  \begin{columns}[T]
    \begin{column}{0.38\textwidth}
      \begin{center}
        \includegraphics[width=0.42\linewidth]{\departmentlogo}\\[0.20cm]
        {\Large\bfseries\color{unadmGreenDark}Filosofia del Derecho}\\
        {\small\color{unadmMuted}\facultyname}
      \end{center}
    \end{column}
    \begin{column}{0.58\textwidth}
      \begin{block}{Datos para copiar a cada actividad}
        \textbf{Alumno:} \studentname\\
        \textbf{Matricula:} \studentid\\
        \textbf{Asignatura:} \coursename\\
        \textbf{Codigo:} \coursecode\\
        \textbf{Figura docente:} \teachingfigure
      \end{block}
      \begin{alertblock}{Archivos base}
        \texttt{\reporttemplate}\\
        \texttt{\coursebib}
      \end{alertblock}
    \end{column}
  \end{columns}
  \accentbar
\end{frame}

\begin{frame}{Lo que ya exige la materia}
  \small
  \begin{tabularx}{\linewidth}{@{}p{0.17\linewidth}Y Y@{}}
    \toprule
    \textbf{Actividad} & \textbf{Producto trabajado} & \textbf{Lectura que debe conservarse} \\
    \midrule
    1 & Mapa conceptual & Filosofia del Derecho como red: objeto, justicia, corrientes, evolucion e importancia practica. \\
    2 & Cuadro comparativo & Derecho Natural y Positivismo como tension entre justicia material y validez institucional. \\
    3 & Glosario + caso & Conceptos juridicos aplicados a matrimonio igualitario y moral constitucional. \\
    4 & Estudio de caso & Poder, Estado y Derecho desde articulos constitucionales y ejemplos periodisticos. \\
    5 & Reflexion + criterios & Interpretacion juridica, hermeneutica y argumentacion en casos penales complejos. \\
    6 & Caso filosofico-juridico & Ayotzinapa como prueba de dignidad, legitimidad, verdad, justicia y limite del poder. \\
    \bottomrule
  \end{tabularx}
  \accentbar
\end{frame}

\begin{frame}{Nucleo narrativo comun}
  \begin{center}
    \begin{tikzpicture}[
      node distance=0.55cm,
      box/.style={
        rounded corners=4pt,
        draw=unadmGreen!40,
        fill=unadmPaper,
        text width=3.20cm,
        minimum height=1.15cm,
        align=center,
        font=\small
      },
      arrow/.style={-{Stealth[length=2.5mm]},very thick,draw=unadmGold}
    ]
      \node[box] (p) {\textbf{Problema}\\Que tension juridica se analiza};
      \node[box,right=of p] (c) {\textbf{Conceptos}\\Que categoria explica el caso};
      \node[box,right=of c] (s) {\textbf{Sistema}\\Que norma, autoridad o institucion opera};
      \node[box,right=of s] (po) {\textbf{Postura}\\Que posicion se defiende};
      \draw[arrow] (p) -- (c);
      \draw[arrow] (c) -- (s);
      \draw[arrow] (s) -- (po);
    \end{tikzpicture}
  \end{center}
  \begin{alertblock}{Uso}
    Esta secuencia sirve para actividades nuevas aunque cambie el producto:
    evita exponer definiciones sueltas y obliga a cerrar con consecuencia.
  \end{alertblock}
  \accentbar
\end{frame}

\unadmsection{Rutas por producto}

\begin{frame}{Matriz de seleccion del producto}
  \scriptsize
  \renewcommand{\arraystretch}{1.18}
  \begin{tabularx}{\linewidth}{@{}p{0.19\linewidth}Y Y Y@{}}
    \toprule
    \textbf{Si la planeacion pide} & \textbf{Estructura de slides} & \textbf{Evidencia visual} & \textbf{Cierre obligatorio} \\
    \midrule
    Mapa conceptual & concepto central, ramas, enlaces, lectura critica & red jerarquica con palabras de enlace & que relacion permite entender el derecho \\
    Cuadro comparativo & criterio, postura A, postura B, sintesis & tabla de 2 o 3 columnas & que enfoque resuelve mejor el problema \\
    Glosario & ubicacion conceptual, definiciones, aplicacion & tarjetas concepto-funcion-caso & que concepto cambia la decision \\
    Estudio de caso & hechos, problema, normas, analisis, postura & matriz hecho-norma-consecuencia & que debe hacer el sistema juridico \\
    Reflexion & tesis, razones, limite, aprendizaje & flujo argumento-consecuencia & que criterio profesional se sostiene \\
    Presentacion oral & tesis, 3 ejes, prueba, postura & slides de apoyo, no texto completo & consecuencia practica y fuente clave \\
    \bottomrule
  \end{tabularx}
  \accentbar
\end{frame}

\begin{frame}{Plantilla: mapa conceptual}
  \begin{columns}[T]
    \begin{column}{0.44\textwidth}
      \begin{block}{Cuando usarla}
        Para actividades como la Semana 2 o cualquier consigna que pida
        identificar relaciones, jerarquias, corrientes, elementos o funciones.
      \end{block}
      \begin{alertblock}{Regla}
        No basta dibujar cajas: cada enlace debe decir que tipo de relacion hay.
      \end{alertblock}
    \end{column}
    \begin{column}{0.52\textwidth}
      \resizebox{0.96\linewidth}{!}{%
      \begin{tikzpicture}[
        every node/.style={font=\scriptsize},
        root/.style={rounded corners=5pt,fill=unadmGreenDark,text=white,align=center,minimum width=2.8cm,minimum height=0.9cm},
        branch/.style={rounded corners=4pt,fill=unadmGreen,text=white,align=center,text width=2.3cm,minimum height=0.70cm},
        leaf/.style={rounded corners=3pt,fill=unadmPaper,draw=unadmGreen!40,align=center,text width=2.25cm,minimum height=0.55cm},
        link/.style={-{Stealth[length=2mm]},draw=unadmGold,thick}
      ]
        \node[root] (r) {Tema juridico};
        \node[branch,above left=0.65cm and 1.00cm of r] (a) {Concepto};
        \node[branch,above right=0.65cm and 1.00cm of r] (b) {Corriente};
        \node[branch,below left=0.65cm and 1.00cm of r] (c) {Norma};
        \node[branch,below right=0.65cm and 1.00cm of r] (d) {Caso};
        \node[leaf,above=0.35cm of a] (a1) {define};
        \node[leaf,above=0.35cm of b] (b1) {contrasta};
        \node[leaf,below=0.35cm of c] (c1) {aplica};
        \node[leaf,below=0.35cm of d] (d1) {prueba};
        \draw[link] (r) -- node[above,sloped]{se explica por} (a);
        \draw[link] (r) -- node[above,sloped]{dialoga con} (b);
        \draw[link] (r) -- node[below,sloped]{se concreta en} (c);
        \draw[link] (r) -- node[below,sloped]{se verifica en} (d);
        \draw[link] (a) -- (a1);
        \draw[link] (b) -- (b1);
        \draw[link] (c) -- (c1);
        \draw[link] (d) -- (d1);
      \end{tikzpicture}
      }
    \end{column}
  \end{columns}
  \accentbar
\end{frame}

\begin{frame}{Plantilla: cuadro comparativo}
  \small
  \begin{tabularx}{\linewidth}{@{}p{0.20\linewidth}Y Y@{}}
    \toprule
    \textbf{Criterio} & \textbf{Enfoque A} & \textbf{Enfoque B} \\
    \midrule
    Fundamento & Que justifica su validez o autoridad. & Que justifica su validez o autoridad. \\
    Funcion juridica & Que problema resuelve dentro del sistema. & Que problema resuelve dentro del sistema. \\
    Riesgo & Que ocurre si se aplica sin limite critico. & Que ocurre si se aplica sin limite critico. \\
    Ejemplo mexicano & Norma, sentencia, institucion o caso. & Norma, sentencia, institucion o caso. \\
    Sintesis & \multicolumn{2}{Y}{Postura: no repetir diferencias; explicar que combinacion o criterio permite resolver mejor el problema.} \\
    \bottomrule
  \end{tabularx}
  \begin{alertblock}{Uso en la materia}
    Replica la logica de Derecho Natural vs. Positivismo: comparar no es poner
    columnas simetricas, sino mostrar una decision entre justicia, validez,
    seguridad juridica, dignidad y eficacia institucional.
  \end{alertblock}
  \accentbar
\end{frame}

\begin{frame}{Plantilla: glosario critico}
  \begin{columns}[T]
    \begin{column}{0.31\textwidth}
      \begin{block}{1. Definicion}
        Concepto breve con fuente: que es y que lugar ocupa en el sistema.
      \end{block}
    \end{column}
    \begin{column}{0.31\textwidth}
      \begin{block}{2. Funcion}
        Que permite decidir, exigir, limitar, validar o interpretar.
      \end{block}
    \end{column}
    \begin{column}{0.31\textwidth}
      \begin{block}{3. Aplicacion}
        Como opera en un caso: matrimonio igualitario, poder estatal, victimas,
        interpretacion penal o desaparicion de personas.
      \end{block}
    \end{column}
  \end{columns}
  \vspace{0.28cm}
  \begin{alertblock}{Formula}
    Definicion base + cita \(\rightarrow\) funcion juridica \(\rightarrow\) consecuencia concreta.
  \end{alertblock}
  \accentbar
\end{frame}

\begin{frame}{Plantilla: estudio de caso}
  \small
  \begin{tabularx}{\linewidth}{@{}p{0.18\linewidth}Y Y@{}}
    \toprule
    \textbf{Paso} & \textbf{Pregunta de trabajo} & \textbf{Salida esperada} \\
    \midrule
    Hechos & Que ocurrio y por que importa juridicamente. & Descripcion sobria, sin dramatizar ni ocultar el dano. \\
    Problema & Que tension se debe resolver. & Pregunta juridica concreta. \\
    Normas & Que articulos, leyes, tesis o principios aplican. & Marco normativo con funcion, no lista decorativa. \\
    Conceptos & Que categorias de Filosofia del Derecho explican el caso. & Poder, Estado, dignidad, validez, moral, interpretacion. \\
    Postura & Que lectura sostengo y que consecuencia produce. & Cierre argumentado y transferible. \\
    \bottomrule
  \end{tabularx}
  \accentbar
\end{frame}

\begin{frame}{Plantilla: interpretacion juridica}
  \begin{columns}[T]
    \begin{column}{0.47\textwidth}
      \begin{block}{Cuatro capas}
        \begin{enumerate}
          \item Texto: que dice literalmente la norma.
          \item Sistema: como se relaciona con otras normas.
          \item Finalidad: que bien juridico protege.
          \item Consecuencia: a quien beneficia o perjudica la lectura.
        \end{enumerate}
      \end{block}
    \end{column}
    \begin{column}{0.49\textwidth}
      \begin{alertblock}{Evitar}
        No presentar la interpretacion como opinion personal aislada. Debe
        apoyarse en hermeneutica, argumentacion, precedentes y efecto real.
      \end{alertblock}
      \begin{block}{Ejemplo trabajado}
        Actividad 5: violencia fisica, incapacidad de resistencia y consentimiento.
      \end{block}
    \end{column}
  \end{columns}
  \accentbar
\end{frame}

\begin{frame}{Plantilla: postura personal}
  \begin{columns}[T]
    \begin{column}{0.48\textwidth}
      \begin{block}{Estructura minima}
        \begin{itemize}
          \item \textbf{Posicion:} sostengo que \ldots
          \item \textbf{Razon:} porque el concepto/norma/caso muestra \ldots
          \item \textbf{Limite:} el riesgo de esta lectura es \ldots
          \item \textbf{Efecto:} para la practica juridica implica \ldots
        \end{itemize}
      \end{block}
    \end{column}
    \begin{column}{0.48\textwidth}
      \begin{alertblock}{Tono}
        Usar primera persona con medida. La postura debe sonar propia, pero no
        coloquial: \emph{Desde esta perspectiva}, \emph{considero que},
        \emph{la posicion que sostengo es}.
      \end{alertblock}
    \end{column}
  \end{columns}
  \accentbar
\end{frame}

\unadmsection{Banco de ejes}

\begin{frame}{Conceptos recurrentes para armar slides}
  \small
  \begin{tabularx}{\linewidth}{@{}p{0.23\linewidth}Y Y@{}}
    \toprule
    \textbf{Eje} & \textbf{Pregunta guia} & \textbf{Uso en actividades} \\
    \midrule
    Validez y eficacia & La norma existe y ademas protege realmente. & Natural/positivismo, matrimonio igualitario, Ayotzinapa. \\
    Derecho y moral & Cuando la moral orienta y cuando excluye. & Moral constitucional, dignidad, normas injustas. \\
    Poder y limite & Quien manda, bajo que reglas y con que control. & Estado mexicano, division de poderes, victimas. \\
    Interpretacion & Que lectura evita formalismo y estereotipos. & Casos penales, jurisprudencia, principios. \\
    Dignidad humana & Que trato exige reconocer a la persona como fin. & Derechos humanos, desaparicion, igualdad, reparacion. \\
    \bottomrule
  \end{tabularx}
  \accentbar
\end{frame}

\begin{frame}{Corrientes y autores como herramientas}
  \begin{columns}[T]
    \begin{column}{0.49\textwidth}
      \begin{block}{No usarlos como lista}
        \begin{itemize}
          \item Garcia Maynez: conceptos juridicos y estructura normativa.
          \item Kelsen: validez, sistema normativo y Estado como orden juridico.
          \item Finnis: bienes humanos, razon practica y justicia.
          \item Laun: relacion y tension entre derecho y moral.
        \end{itemize}
      \end{block}
    \end{column}
    \begin{column}{0.47\textwidth}
      \begin{block}{Usarlos con funcion}
        \begin{itemize}
          \item Definir una categoria.
          \item Contrastar una corriente.
          \item Justificar un criterio.
          \item Evaluar una norma o sentencia.
          \item Cerrar una postura personal.
        \end{itemize}
      \end{block}
    \end{column}
  \end{columns}
  \accentbar
\end{frame}

\begin{frame}{Sistema mexicano: prueba de realidad}
  \begin{columns}[T]
    \begin{column}{0.32\textwidth}
      \begin{block}{Normas}
        Constitucion, Ley de Amparo, Ley General de Victimas, leyes especiales,
        codigos penales y legislacion secundaria.
      \end{block}
    \end{column}
    \begin{column}{0.32\textwidth}
      \begin{block}{Instituciones}
        SCJN, tribunales, autoridades administrativas, fiscalias, congreso,
        organos electorales y organismos de derechos humanos.
      \end{block}
    \end{column}
    \begin{column}{0.32\textwidth}
      \begin{block}{Casos}
        Matrimonio igualitario, poder estatal, interpretacion penal, victimas,
        desaparicion de personas y control constitucional.
      \end{block}
    \end{column}
  \end{columns}
  \accentbar
\end{frame}

\begin{frame}{Fuentes: criterio de uso en la exposicion}
  \small
  \begin{tabularx}{\linewidth}{@{}p{0.24\linewidth}Y Y@{}}
    \toprule
    \textbf{Tipo de fuente} & \textbf{Funcion en slides} & \textbf{Error a evitar} \\
    \midrule
    Bibliografia teorica & Define, distingue o fundamenta conceptos. & Mencionar autores sin explicar que aportan. \\
    Norma vigente & Muestra el marco obligatorio de decision. & Copiar articulos largos sin interpretarlos. \\
    Jurisprudencia & Prueba como se aplica un concepto en conflicto real. & Reducirla a resumen de expediente. \\
    Nota periodistica & Aterriza el problema en un hecho verificable. & Convertir la exposicion en narracion de noticias. \\
    Reporte propio & Da tesis, postura, estructura y cierre. & Leer parrafos completos en pantalla. \\
    \bottomrule
  \end{tabularx}
  \accentbar
\end{frame}

\unadmsection{Actividades modelo}

\begin{frame}{Actividades 1--3: de conceptos a aplicacion}
  \small
  \begin{tabularx}{\linewidth}{@{}p{0.20\linewidth}Y Y@{}}
    \toprule
    \textbf{Actividad} & \textbf{Claim de presentacion} & \textbf{Slide clave} \\
    \midrule
    A1: mapa conceptual & La Filosofia del Derecho ordena el fenomeno juridico porque conecta justicia, norma, corrientes y practica. & Mapa con jerarquias y palabras de enlace. \\
    A2: cuadro comparativo & El jurista necesita legalidad positiva y critica material de justicia; una sin la otra produce formalismo o abstraccion. & Tabla comparativa con sintesis final. \\
    A3: glosario y caso & Los conceptos juridicos no son definiciones escolares: permiten diagnosticar discriminacion, validez y eficacia. & Tarjetas de concepto + matriz de caso. \\
    \bottomrule
  \end{tabularx}
  \accentbar
\end{frame}

\begin{frame}{Actividades 4--6: del poder al caso limite}
  \small
  \begin{tabularx}{\linewidth}{@{}p{0.20\linewidth}Y Y@{}}
    \toprule
    \textbf{Actividad} & \textbf{Claim de presentacion} & \textbf{Slide clave} \\
    \midrule
    A4: Estado y Derecho & El Estado de Derecho existe cuando poder politico, social y juridico quedan sometidos a limites constitucionales. & Matriz articulo-concepto-ejemplo. \\
    A5: interpretacion & Interpretar decide que derechos protege realmente el sistema y que sesgos corrige o reproduce. & Flujo texto-sistema-finalidad-consecuencia. \\
    A6: Ayotzinapa & La desaparicion forzada prueba si el derecho funciona como garantia de dignidad o solo como legalidad declarativa. & Caso: hecho-problema-concepto-norma-postura. \\
    \bottomrule
  \end{tabularx}
  \accentbar
\end{frame}

\begin{frame}{Slide base: caso filosofico-juridico}
  \begin{columns}[T]
    \begin{column}{0.49\textwidth}
      \begin{block}{Caso}
        \textit{Editable:} identifica el hecho o criterio judicial que pone a
        prueba una categoria de la Filosofia del Derecho.
      \end{block}
      \begin{block}{Problema}
        \textit{Editable:} formula una pregunta: que revela el caso sobre
        validez, justicia, moral, poder, dignidad o interpretacion.
      \end{block}
    \end{column}
    \begin{column}{0.47\textwidth}
      \begin{block}{Aplicacion}
        \textit{Editable:} vincula conceptos, normas y fuentes con una
        consecuencia juridica especifica.
      \end{block}
      \begin{alertblock}{Postura}
        \textit{Editable:} sostengo que \ldots porque \ldots y esto implica \ldots
      \end{alertblock}
    \end{column}
  \end{columns}
  \accentbar
\end{frame}

\unadmsection{Control de calidad}

\begin{frame}{Checklist antes de entregar}
  \begin{itemize}
    \item La portada coincide con actividad, semana, materia, docente y fecha.
    \item La presentacion responde a la planeacion sin copiar instrucciones.
    \item Cada slide tiene funcion clara: problema, concepto, evidencia o cierre.
    \item El producto visual pedido aparece como evidencia principal.
    \item Hay postura personal argumentada, no solo resumen de fuentes.
    \item Las fuentes citadas en el reporte sostienen tambien la exposicion.
    \item El cierre traduce el tema a una consecuencia juridica concreta.
  \end{itemize}
  \accentbar
\end{frame}

\begin{frame}{Cierre editable}
  \begin{block}{Hallazgo principal}
    \textit{Editable:} La actividad permite concluir que \ldots
  \end{block}
  \begin{block}{Criterio profesional}
    \textit{Editable:} Para la practica juridica, este tema exige \ldots
  \end{block}
  \begin{alertblock}{Ultima frase}
    No cerrar con ``eso fue todo''. Cerrar con una idea defendible:
    norma, justicia, poder y persona deben leerse juntos.
  \end{alertblock}
  \accentbar
\end{frame}

\end{document}
